% SST-REPORT-SECTION:IDENTIFICATION
% SST-REPORT-SECTION:QUESTION
% SST-REPORT-SECTION:HYPOTHESES
% SST-REPORT-SECTION:ASSUMPTIONS
% SST-REPORT-SECTION:SYMBOLS
% SST-REPORT-SECTION:EQUATIONS
% SST-REPORT-SECTION:ALGORITHM
% SST-REPORT-SECTION:SOURCES
% SST-REPORT-SECTION:GATES
% SST-REPORT-SECTION:NUMERICS
% SST-REPORT-SECTION:BACKENDS
% SST-REPORT-SECTION:BLINDNESS
% SST-REPORT-SECTION:STATISTICS
% SST-REPORT-SECTION:RESULTS
% SST-REPORT-SECTION:INTERPRETATION
% SST-REPORT-SECTION:REPRODUCIBILITY
\documentclass[11pt,a4paper]{article}
\usepackage[margin=22mm]{geometry}
\usepackage[T1]{fontenc}
\usepackage[utf8]{inputenc}
\usepackage{amsmath,amssymb,bm,booktabs,longtable,hyperref,enumitem}
\newcommand{\vswirl}{\mathbf{v}_{\!\boldsymbol{\circlearrowleft}}}
\newcommand{\rhoF}{\rho_{\!f}}
\newcommand{\rhoE}{\rho_{\!E}}
\newcommand{\rhoM}{\rho_{\!m}}
\newcommand{\rc}{r_c}
\title{A062---SST Microgravity Force-Gradient Falsifier\\\large v0.2.0: E011 3\_1 Finite-Core Kelvin Impulse Interaction}
\author{SST-Workbench / scientific implementation note}
\date{12 October 2026}
\begin{document}\maketitle
\section{Identification and limits}
This is the protocol of a restricted, *instantaneous* two-vortex filament experiment using actual E011 static-ready geometry. A separate standalone PDF is generated by \texttt{run\_science.cmd FULL}. The framework-owned result insert follows when a valid local frozen v1.0.7 framework is available. A publication benchmark that has already been inspected is not independent blind validation. Static-ready does not mean dynamics-ready.
\IfFileExists{AUTO_RESULTS.tex}{\input{AUTO_RESULTS.tex}}{}
\section{Research question and hypotheses}
Can canonical-circulation regularized vortex filaments with independent source-qualified trefoil centrelines yield an orientation-independent attractive $R^{-2}$ interaction, without inserting $G$? The null is that at least one of the far-field exponent, universal attraction, or cross-provider amplitude gates fails; the restricted candidate alternative requires all to pass. These thresholds are model tests, not general SST falsification criteria.
\section{Physical and geometric assumptions}
The background is 3D Euclidean, unforced, inviscid, and incompressible. Knot centrelines are fixed *initial geometries*; no Euler equilibrium or stability is assumed. The velocity field is regularized by the Rosenhead core kernel. Both circulation magnitudes use the SST canon. Mapping geometric reach to the canonical hydrodynamic core radius is an additional model assumption; it is not derived. No mass of a trefoil is assumed, and Newton's constant does not enter the candidate solver.
\section{Units, symbols, and equations}
Let $\Gamma$ have units $\mathrm{m^2/s}$, $\rhoF$ have $\mathrm{kg/m^3}$, $\mathbf{X}$ have $\mathrm m$, $\mathbf{I}$ have $\mathrm{kg\,m/s}$, $\mathbf{F}$ have $\mathrm N$, and $K$ have $\mathrm{N/m}$. The canonical parameters are:
\[\rhoF=7.0\times 10^{-7}\,\mathrm{kg/m^3},\quad \rc=1.40897017\times10^{-15}\,\mathrm m,\quad |\vswirl|=1.09384563\times10^6\,\mathrm{m/s}.\]
\begin{align}
\Gamma&=2\pi\rc|\vswirl|,\\
\mathbf v_A(\mathbf x)&=\frac{\Gamma_A}{4\pi}\oint_A\frac{d\mathbf X_A\times(\mathbf x-\mathbf X_A)}{(|\mathbf x-\mathbf X_A|^2+\rc^2)^{3/2}},\\
\mathbf I_B&=\frac{\rhoF\Gamma_B}{2}\oint_B\mathbf X_B\times d\mathbf X_B,\\
\left.\frac{d\mathbf I_B}{dt}\right|_{A\to B}&=\rhoF\Gamma_B\oint_B\mathbf v_A(\mathbf X_B)\times d\mathbf X_B,\\
K(R)&=\frac{dF_{B,z}}{dR},\qquad n=-\frac{d\log|F_{B,z}|}{d\log R}.
\end{align}
$[\rhoF\Gamma_B v_A\,dX]=\mathrm N$. Here $d\mathbf I_B/dt$ measures exchange of filament Kelvin impulse; identifying it with a physical force on an isolated material particle is an additional unproved assumption. The total fluid momentum cannot be inferred from the filament impulse alone.
\subsection{Pressure-Poisson consistency diagnostic}
\begin{align}
\nabla^2p&=-\rhoF\partial_i v_j\partial_j v_i,\\
\nabla^2 p_{\times}&=-\rhoF(\partial_i v_{A,j}\partial_j v_{B,i}+\partial_i v_{B,j}\partial_j v_{A,i}).
\end{align}
For a periodic diagnostic grid, the zero Fourier mode is subtracted and $\widehat p_{\times}(\mathbf k)=-\widehat S_{\times}(\mathbf k)/|\mathbf k|^2$ for $\mathbf k\neq0$. Small-grid algebraic residual is \emph{not} a resolved free-space pressure or force. The canonical core remains unresolved at the selected CPU grid sizes.
\section{Exact method and source-independent numerics}
Input geometric centrelines are actual E011 v0.3.0 \texttt{3\_1} Gilbert and KnotPlot provider anchors. The pinned originals are a Gilbert \texttt{Ideal.txt.gz} Fourier record (\texttt{3:1:1}) and a KnotPlot closed Geomview \texttt{VECT}. Both source SHA-256 values are validated against E011's source locators. Geometry is resampled at constant arc length to $N=48,96,192,384$ segments and normalized to perimeter one. A proper rotation aligns the area vector with $+z$; reflections are \emph{not} silently applied, so opposite provider signed-writhe conventions remain visible.

Set $L_{\mathrm{phys}}=\rc/\mathrm{reach}_{\mathrm{E011}}$ separately for each upstream provider. Evaluate the regularized velocity at each test filament segment midpoint. Multiply the dimensionless Kelvin exchange by $\rhoF\Gamma_A\Gamma_B$. Sweep $R/L=1,2,3,5,8,12,18,27,40$ and test relative tilts $0,\pi/2,\pi$. Fit the logarithmic slope only at $R/L\ge8$. Compute a central difference at $R/L=5$, keeping the $1/L_{\mathrm{phys}}$ gradient conversion explicit. The two provider families count as two geometry providers, not two experiments; different $N$ are not independent samples.
\section{Numerical gates and stopping rules}
Scientific gates never short circuit after a physical fail. Source hashes, geometry closure, Kelvin impulse reciprocity (relative residual $<2\times10^{-3}$), quadrature refinement (relative change $<0.025$), $|n-2|\le0.2$ for both providers, sign persistence between $0$ and $\pi$ tilt, and two-provider amplitude agreement ($\le5\%$ at $R/L=5$) are scored independently. Real C++ FP64 parity is required for its own numerical gate ($<2\times10^{-8}$ relative vector error); unavailable native hardware is UNRESOLVED, never a Python fallback. No SYCL results are asserted. The periodic Poisson grid is explicitly unresolved physically; no 3D Euler temporal trajectory, Kelvin/Floquet stability, isolated momentum-flux pressure closure, derived mass or gravitational constant follows.
\section{Blindness, uncertainty and interpretation}
This research protocol was constructed after the article and earlier Gaussian results were public. It is an exploratory model comparison, \emph{not} a genuine blinded held-out force measurement. The supplied sources are original upstream geometry data and their immutable E011 static qualification summaries. File hashes and source contracts distinguish metadata and source geometry; later model parameters must not be retuned without a new version. There is no claim to independently extract $G$ or any SST canon constant from the published force curve. The original article concerns gram-sized apparatus, not isolated elementary vortex knots.
\section{Results and reproducibility}
The standalone science output is always created by \texttt{run\_science.cmd FULL} after \texttt{FREEZE}; it includes \texttt{SCIENCE\_RESULTS.json}, \texttt{GATE\_LEDGER.json}, figures and a PDF. The shared framework v1.0.7 is pinned by the generated \texttt{.sst\_framework\_root} file. \texttt{run\_all.cmd} requires the genuine external frozen framework and may fail closed if it is unavailable. Source, report, gates, science contract, and raw E011 input hashes are never modified after preregistration. A failing physical model must not be rescued by changing geometry or thresholds after examining the results.
\begin{thebibliography}{9}
\bibitem{Saffman1992}P. G. Saffman (1992), \emph{Vortex Dynamics}, Cambridge University Press. \url{https://doi.org/10.1017/CBO9780511624063}.
\bibitem{Batchelor1967}G. K. Batchelor (1967), \emph{An Introduction to Fluid Dynamics}, Cambridge University Press. \url{https://doi.org/10.1017/CBO9780511800955}.
\bibitem{Zhu2026}Y. Zhu et al. (2026), ``Stopped-light-enhanced gravitational force sensing,'' \emph{Nature Nanotechnology}. \url{https://doi.org/10.1038/s41565-026-02298-8}.
\bibitem{E011}SST-Workbench (2026), E011/SKLSA v0.3.0, \emph{Static Seed Atlas}, upstream E010/PKLSA v0.3.1; internal source and provider qualification contracts.
\end{thebibliography}
\end{document}
